\section{Result Records and Comparison Rules}\label{sec:scope}
The unit of comparison is one evaluated condition, rather than a paper or a headline score. A result record binds a source version and locator to its target, receiver, intervention, attacker procedure, workload and outcomes (Table~\ref{tab:result-record}). A family can contribute several records when it reports different execution modes, protected views or cost components. This is why a valid result can answer one deployment question while leaving another unresolved.
\begin{table}[t]\centering\footnotesize
\caption{A single result record. Only fields relevant to the requested comparison must align; unavailable required fields remain unresolved.}\label{tab:result-record}
\begin{tabularx}{\columnwidth}{@{}p{.24\columnwidth}Y@{}}\toprule
Record field & Required meaning\\\midrule
Source & Family/result ID, version, source locator, reading scope and hash.\\
Target & Hidden property, inference goal and protection relation: confidentiality, attribution, human/agent distinction or authorized diagnosis.\\
Receiver & Actual permission and released representation, including session/document linkage and any auxiliary inputs.\\
Intervention & Where protection ends; whether edits affect stored traces, transmission or executed actions.\\
Attacker & Fit/calibration condition, test population, training access and adaptation budget.\\
Workload & Task, system, topology/hardware and grouping/holdout.\\
Outcomes & Inference metric/denominator; objective or proxy utility; applicable cost object; uncertainty and field states.\\\bottomrule
\end{tabularx}\end{table}


\paragraph{One record, two indexing aids}
\scope{} indexes the record's Secret/target, Channel, Observer, Pipeline and Exploitation goal. The Evidence Plane is the record's source and evaluation columns, not a second taxonomy or an evidence score. Pipeline indices distinguish hardware (P1), serving (P2), transport (P3), application/API (P4), orchestration (P5), tools/MCP (P6), environment (P7) and multi-agent communication (P8). A generating layer and a receiving layer can differ. Permission, activity, calibration access and protection relation are separate fields: an authorized fault-localization component is not thereby a confidentiality attack, and an attribution label is not automatically a protected secret.

\paragraph{Three comparison outcomes}
A \emph{compatible} comparison has evidence for the fields needed by the requested inference. An \emph{incompatible} comparison has a positively established mismatch, such as token payload versus wire bytes or stored timestamps versus executed actions. An \emph{unresolved} comparison lacks a necessary binding or transfer test. Only the affected inference is withheld: an unreported task endpoint does not invalidate a measured bandwidth cost. Table~\ref{tab:comparison-contracts} applies these outcomes to descriptive comparison, paired effects and deployment transfer.
\begin{table}[t]\centering\footnotesize
\caption{Comparison rules. A known mismatch and a missing binding have different consequences.}\label{tab:comparison-contracts}
\begin{tabularx}{\columnwidth}{@{}p{.28\columnwidth}Y@{}}\toprule
Requirement & Condition and resulting license\\\midrule
Descriptive comparison & Align target, receiver permission, metric meaning and cost object. Different systems may be described together; token payload and wire bytes cannot be ranked as the same cost.\\
Paired effect & Compare the same defended test population and outcome. A known population mismatch makes that subtraction incompatible; an unverified fit binding leaves the effect unresolved.\\
Deployment transfer & Establish a bridge for the changed receiver, intervention, workload, topology or hardware. A demonstrated violation blocks the proposed transfer; an untested bridge leaves it unresolved. Permit only the range and endpoints actually covered.\\\bottomrule
\end{tabularx}\end{table}


\paragraph{Missingness and adaptation}
\NRR{} records a quantity unreported within the inspected source scope; \NRA{} records unavailable material; \NRV{} records an unverified field, version or run binding; \NA{} records inapplicability. A missing field gives an unresolved judgment, not an observed mismatch. Fitting on protected data (D1) and comparing clean- versus protected-trained procedures on the same defended test population (D2) are separate entries in the attacker column. Neither alone establishes that the observer can obtain the labeled data or afford the fitting budget.

\paragraph{A complete traceable example}
The chain below connects a family, a result, the inspected source and the deployment judgment. The same identifiers join the supplied family index, result ledger and claim links; full source hashes and reading scopes remain in the companion evidence package.
\begin{table}[t]\centering\footnotesize
\caption{Filled record and its licensed judgment. Values are source-reported; no cross-mode effect is computed.}\label{tab:result-lineage}
\begin{tabularx}{\columnwidth}{@{}p{.24\columnwidth}Y@{}}\toprule
Join key / field & Entry\\\midrule
Family / result & \texttt{dangol2026workflow} / \texttt{DANGOL-livebinding}.\\
Version / locator & arXiv:2606.07150v3, \S9.7, OHTTP paragraph after Table~5 \cite{dangol2026workflow}.\\
Target / view & Workflow classification from the evaluated live OHTTP/A2A traffic view; not GUI-state observation.\\
Intervention & Live fixed-rate cover with fresh per-session identifiers, distinct from transformed stored captures.\\
Workload / testbed & Workflow classes with client, relay, gateway and agent co-located on one loopback host.\\
Outcome / cost & Accuracy 0.30; about 13$\times$ bytes and 11$\times$ latency in that condition.\\
Unresolved fields & Objective terminal task utility \NRR{}; precise paired live-refit comparison \NRV{}.\\
Licensed judgment & These costs describe this live protocol condition. They neither price dummy GUI actions nor prove task preservation. A matched objective endpoint is needed for the latter.\\\bottomrule
\end{tabularx}\end{table}


\paragraph{Scope of the method claim}
A complete structured checklist can express these same checks. The separate model-assisted comparison favors that checklist on its original joint endpoint (Appendix~\ref{app:scope-workflow-boundary}). We use \scope{} to organize the resulting knowledge; improved review accuracy, efficiency and human usability are unestablished.
